\section{Evaluation}\label{sec:eval}

This section reports artifact-backed measurements for xlog's deterministic, probabilistic, and neural-symbolic subsystems. Every quantitative claim is tied to a checked-in benchmark artifact. \emph{These are single-system measurements on development hardware: ablations of xlog against xlog's own baselines, not controlled head-to-head comparisons against other engines} (\Cref{sec:limits}). Throughput-target envelopes for operations measured only in isolation are maintained as an internal regression contract and are not restated here.

\subsection{Methodology}

All measurements were collected on an NVIDIA RTX PRO 3000 (Blackwell, 12\,GB VRAM, compute capability 12.0). Rust crates are built in release mode; CUDA modules are staged cubin-first with PTX fallback and explicit JIT warm-up; Python components use \texttt{pyxlog} with CUDA-enabled PyTorch. The benchmark harness uses Criterion.rs with the settings in \Cref{tab:stat-protocol}, and GPU measurements include warm-up to amortize kernel-module loading and CUDA context setup.

\begin{table}[htbp]
  \centering
  \caption{Benchmark harness settings.}
  \label{tab:stat-protocol}
  \small
  \begin{tabular}{@{}ll@{}}
    \toprule
    \textbf{Setting} & \textbf{Value} \\
    \midrule
    Sample size       & 10--100 runs (benchmark-dependent) \\
    Warm-up           & 3 iterations (PTX JIT, memory-pool init) \\
    Significance level & 0.10 \\
    Noise threshold   & 0.05 (ignore ${<}5\%$ variance) \\
    Seeding           & Deterministic LCG \\
    \bottomrule
  \end{tabular}
\end{table}

\subsection{Worst-Case Optimal Joins}

The WCOJ subsystem (\Cref{sec:datalog}) is measured against the binary-join baseline on four paper-scale fixtures---call-graph edges, Andersen points-to~\cite{andersen1994}, \texttt{ddisasm}-style disassembly~\cite{ddisasm}, and a recursive neural-symbolic mining analog---each with $4.19$M input rows. We measure at this scale deliberately: WCOJ speedups depend strongly on input size and skew, and synthetic micro-benchmarks systematically overstate them, so multi-rule recursive workloads at millions of tuples are the honest setting. Both paths must produce identical row sets before timing, and every measurement is a median over 10 runs with coefficient of variation ${\le}0.05$.

\Cref{fig:wcoj} reports the speedups: a $27.96\times$ geometric mean, with every fixture between $26.6\times$ and $29.6\times$.

\begin{figure}[htbp]
  \centering
  \includegraphics{results.pdf}
  \caption{WCOJ vs.\ binary-join speedup on four paper-scale fixtures ($4.19$M input rows each). Single-system ablation; the dashed line is the geometric mean.}
  \label{fig:wcoj}
\end{figure}

On uniform or empty inputs the adaptive classifier falls back to the binary-join chain. Peak observed allocation was about $2.3$\,GB, within the device's 12\,GB budget.

\subsection{Circuit Caching and Training}

Neural-symbolic training is measured on MNIST single-digit addition (512 training images, 5 epochs, batch size 64). Circuit caching (\Cref{sec:prob}) amortizes D4 compilation across epochs: a cache ablation over 3 seeds and 3 epochs reduced mean training time from $242.9$\,s (uncached) to $88.9$\,s (cached), a $2.74\times$ speedup (95\% CI $[2.29, 3.18]$). \Cref{tab:training-summary} summarizes the full-run timing: the first epoch pays cold-start compilation and verification, while warm epochs run only the cached forward/backward kernels. The ablation is a separate experiment from this full run---its uncached arm recompiles the circuit every epoch while its cached arm compiles once, so it isolates compilation amortization; the apparent speedup therefore grows with epoch count toward the cold/warm ratio rather than being a fixed property.

\begin{table}[htbp]
  \centering
  \caption{MNIST-addition training summary (mean over 3 seeds).}
  \label{tab:training-summary}
  \small
  \begin{tabularx}{\linewidth}{@{}l l X@{}}
    \toprule
    \textbf{Metric} & \textbf{Value} & \textbf{Notes} \\
    \midrule
    First epoch (cold)         & $71.7$\,s  & D4 compile + CDCL verify \\
    Steady-state epoch (warm)  & $0.25$\,s  & cached forward/backward path \\
    Total training (5 epochs)  & $72.6$\,s  & \\
    Per-query cost             & $0.97$\,ms & forward + backward \\
    Circuit-cache speedup      & $2.74\times$ & 95\% CI $[2.29, 3.18]$ \\
    \bottomrule
  \end{tabularx}
\end{table}

\subsection{Runtime Optimization}

The runtime optimizations (\Cref{sec:datalog}) are measured on repeated-session fixtures. The persistent hash-index manager records a $3.21\times$ speedup on a build-heavy repeated-session semi-join (cached median $0.08$\,s vs.\ uncached $0.26$\,s) with zero tracked host transfers, and a profile-gated shared-memory chain scorer records a $5.58\times$ speedup on a chain-hot fixture with a bounded $1.3\%$ regression on the small-input control. Common-subexpression elimination and adaptive re-optimization are validated for output equivalence rather than reported as headline speedups. These two figures are point medians on single fixtures; per-fixture dispersion is left to future work.

\subsection{Capability Comparison}

\Cref{tab:comparison} positions xlog qualitatively against representative systems; each addresses a subset of the design space, and xlog's contribution is unifying them in one GPU-resident runtime. The comparison is qualitative: xlog has not been benchmarked head-to-head against these systems on identical workloads.

\begin{table*}[htbp]
  \centering
  \caption{Capability comparison across neurosymbolic and GPU logic systems.}
  \label{tab:comparison}
  \small
  \begin{tabularx}{\linewidth}{@{}l X X X X X@{}}
    \toprule
    \textbf{Dimension} & \textbf{DeepProbLog} & \textbf{Scallop} & \textbf{Lobster} & \textbf{GPUlog} & \textbf{xlog} \\
    \midrule
    Symbolic execution      & CPU (Prolog) & CPU (Datalog)    & GPU (Datalog)      & GPU (HISA) & GPU (semi-naive) \\
    Probabilistic inference & Exact (CPU)  & Provenance (CPU) & Provenance (GPU)   & None       & Exact d-DNNF + MC (GPU) \\
    Circuit verification    & None         & None             & None               & ---        & On-GPU CDCL \\
    Neural integration      & Prolog bridge & Differentiable  & Differentiable (GPU) & None     & Zero-copy, fused \\
    Zero-copy ML interop    & No           & Partial          & Partial            & No         & Yes \\
    Epistemic reasoning     & No           & No               & No                 & No         & Yes (finite fragment) \\
    \bottomrule
  \end{tabularx}
\end{table*}

\subsection{CUDA Certification}

The CUDA kernel certification suite covers all GPU kernel operations: 207 tests across 33 categories (25 infrastructure C01--C25 + 8 GPU-specific G01--G08), at 100\% pass on the latest full certification, covering 22 kernel modules. The infrastructure categories cover toolchain validation, launch configuration, memory spaces, synchronization, warp-level operations, atomics, floating-point and determinism semantics, and host-device transfer accounting; the GPU-specific categories validate circuit forward/backward evaluation, weight injection, transfer efficiency, circuit caching, PTX robustness, and SAT/CDCL solving. Because the public repository does not rely on hosted GPUs, the suite runs through GPU-backed local release validation; the always-on CI covers formatting, linting, and no-GPU build checks.
